% =========================
% Cas pratique 11
% =========================

\section[Cas 11]{Cas pratique 11 : Médiation judiciaire}

%--------------------------
\begin{frame}{Question 1}

\begin{longtblr}{colspec = {X[1,l] X[2,l] X[2,l]}, rowsep = 3pt}

Art. \textbf{750-1} CPC
& \textbf{Avantages}
& \textbf{Inconvénients} \\

\hline
\textbf{Conciliation} (art. \textbf{21} CPC)
&
Très rapide \newline
Faible coût \newline
Adaptée aux litiges simples
&
Peu efficace en contentieux technique complexe \newline
Accord souvent difficile à obtenir
\\

\textbf{Médiation judiciaire} (art. \textbf{131-1} et s.)
&
Confidentialité \newline
Souplesse des échanges \newline
Utile pour négocier licence/transaction
&
Coût du médiateur \newline
Dépend de la coopération des parties
\\

\textbf{Procédure participative}
(art. \textbf{2062} s. C.civ. ; art. \textbf{1542} s. CPC)
&
Encadrée par les avocats \newline
Travail technique structuré \newline
Peut accélérer la résolution
&
Nécessite forte collaboration \newline
Peu utilisée en pratique
\\

\end{longtblr}

\end{frame}

%--------------------------
\begin{frame}{Question 2}

\begin{block}{}
Loi n°95-125 du 8 février 1995 (médiation judiciaire).
+ Art. \textbf{131-1} à \textbf{131-15} CPC.
+ art. \textbf{131-14} CPC et art. \textbf{21-3} CPC
(confidentialité)
+ art. \textbf{131-9} CPC (information au juge).
\end{block}

\begin{itemize}
  \item \textbf{Entretiens séparés :}
  le médiateur peut organiser des caucus pour recevoir des informations
  qu’une partie ne souhaite pas révéler immédiatement.

  \item \textbf{Secret et loyauté :}
  le médiateur est tenu à la confidentialité et les parties doivent participer
  de bonne foi.

  \item \textbf{Risques si échec :}
  les échanges restent en principe inutilisables au procès,
  mais l’information peut être découverte par d’autres voies
  (expertise, communication de pièces), avec un risque stratégique pour FUNITALIA.
\end{itemize}

\end{frame}

%--------------------------
\begin{frame}{Question 3}

\begin{longtblr}{colspec = {X[1,l] X[1,l]}, rowsep = 1pt}

\textbf{DUTCHJETSKI} & \textbf{FUNITALIA} \\

\hline

Réduire l’aléa judiciaire (nullité + débat technique complexe). 
& Éviter une interdiction et des dommages-intérêts élevés. \\

Obtenir une solution globale (France/Italie + ventes Espagne). 
& Régler confidentiellement les ventes sous autre appellation. \\

Limiter les coûts et la durée (experts, saisies, contentieux long). 
& Négocier une licence, un calendrier ou un écoulement de stock. \\

Accéder à des informations commerciales via accord amiable. 
& Réduire les risques financiers et réputationnels. \\

\end{longtblr}

\end{frame}